\section{Introducción}

Los Transformers resuelven los problemas de varios pasos aplicando una sucesión
de capas, cada una con sus propios parámetros. La profundidad parece clave para
el razonamiento \citep{saunshi2025reasoning}, pero añadir capas encarece de forma
lineal los parámetros, la memoria y el entrenamiento. El Universal
Transformer \citep{dehghani2019universal} ya aplica de forma recurrente un mismo
bloque; trabajos recientes lo estudian como \emph{Looped Transformer}: un único
bloque aplicado $N$ veces sobre su propia salida, con profundidad de cómputo
creciente y parámetros constantes
\citep{giannou2023looped,saunshi2025reasoning,geiping2025scaling}.

Estudiamos esta propuesta a pequeña escala y con control total de los datos:
entrenamos desde cero GPT y Looped Transformers de pocos millones de parámetros
en problemas aritméticos en castellano cuya dificultad es el número de
operaciones. Buscamos responder a \textbf{(P1)}~si un bucle de $N$ iteraciones
iguala a un GPT de $N$ capas con hasta $N$ veces menos parámetros (al compartir el bloque),
\textbf{(P2)}~si las iteraciones necesarias crecen con la dificultad y
\textbf{(P3)}~si el modelo generaliza a problemas con más operaciones que los
vistos al entrenar.\ifinterno{} Además, \textbf{(P4)}~analizamos el interior del bucle con \emph{logit lens} y métricas de representación y atención.\fi{}

Aportamos una implementación propia, sobre el código de \citet{raschka2024build},
de ambos modelos con una interfaz común y validada numéricamente
(Sec.~\ref{sec:validacion}); un conjunto de problemas sintéticos con seis niveles
y los dos más difíciles reservados para evaluación (Sec.~\ref{sec:datos}); y una
comparación controlada a igual cómputo y a igual número de parámetros, con
barrido de iteraciones y análisis interno del bucle
(Secs.~\ref{sec:experimentos} y~\ref{sec:analisis}).
